%% ma: L12-B19-0016
%% lop: 12
%% chuong: 6
%% bai: 19
%% dang: 12.19.1
%% loai: TLN
%% muc: VD
%% dap_an: "58"
%% nguon: "(NBV)-ĐỀ SỐ 10-MỖI NGÀY 1 ĐỀ THI 2025.pdf (D10, MỖI NGÀY 1 ĐỀ - Đề số 10), Phần III Câu 6"
%% kien_thuc: "Bài 19 (công thức xác suất toàn phần và công thức Bayes)"
%% trang_thai: chinh-thuc
%% ghi_chu: "Cùng ý tưởng với L12-B19-0013 và L12-B19-0011 (xác suất toàn phần và Bayes từ xác suất có điều kiện hai bước), nhập vào dạng 12.19.1 đã duyệt. Hình gốc là ảnh học sinh trước máy tính, trang trí."
\begin{ex}
Trong kì thi tốt nghiệp trung học phổ thông, trường THPT A có $60\%$ học sinh lựa chọn khối D để xét tuyển đại học. Biết rằng, nếu một học sinh lựa chọn khối D thì xác suất để học sinh đó đỗ đại học là $0{,}7$ còn nếu học sinh không lựa chọn khối D thì xác suất để học sinh đó đỗ đại học là $0{,}8$. Chọn ngẫu nhiên một học sinh của trường THPT A đã tốt nghiệp trong kì thi trên. Giả sử xác suất để học sinh đó đã chọn khối D biết học sinh này đã đỗ đại học là $\dfrac mn$ với $n$ là số nguyên dương và $\dfrac mn$ là phân số tối giản. Tính giá trị của $m+n$.
\shortans{58}
\loigiai{Gọi $A$: ``học sinh chọn khối D'', $B$: ``học sinh đỗ đại học''. Ta có $P(A)=0{,}6$, $P(\overline A)=0{,}4$, $P(B\mid A)=0{,}7$, $P(B\mid\overline A)=0{,}8$. Theo công thức xác suất toàn phần, $P(B)=0{,}6\cdot0{,}7+0{,}4\cdot0{,}8=0{,}74$. Theo công thức Bayes,
\[P(A\mid B)=\dfrac{P(A)P(B\mid A)}{P(B)}=\dfrac{0{,}42}{0{,}74}=\dfrac{21}{37}.\]
Vậy $m=21$, $n=37$ và $m+n=58$.}
\end{ex}
